
\subsubsection{Findings}

\textbf{Finding 1: Researchers engage with AI documentation assistance.} The 92\% acceptance rate establishes that the engagement mechanism works—researchers will adopt the tool when it reduces friction. However, this validates user engagement (will they use it?), not downstream quality improvement (does use lead to better documentation?). The latter requires evaluation beyond proof-of-concept scope.

\textbf{Finding 2: Cross-domain generalization emerges without domain-specific tuning.} Minimal variance across vision (89.7\%), NLP (89.8\%), and tabular (88.8\%) datasets indicates that contextual prompting with domain-stratified exemplars captures cross-cutting documentation needs. A single model serving all domains reduces deployment complexity.

\textbf{Finding 3: High modification rates indicate genuine utility.} The 26.8\% modification rate demonstrates users invest effort to refine suggestions they find helpful. If suggestions were low-quality, users would either reject them (write from scratch) or show minimal engagement (accept without reading). The observed pattern suggests suggestions provide substantive value as starting points.

\textbf{Finding 4: The proof-of-concept establishes necessary but not sufficient conditions.} High acceptance validates that researchers will engage with AI suggestions—the critical prerequisite for downstream quality gains. However, whether engagement translates to improved documentation completeness, semantic quality, or reduced time burden remains unmeasured.

\subsubsection{Limitations}

\textbf{Limitation 1: Proof-of-concept validated acceptance but not downstream outcomes.} We measured acceptance rate (92\%) but did not evaluate whether high acceptance translates to improved documentation completeness, expert-rated quality, or reduced time-to-completion. The causal chain from "users accept suggestions" to "documentation quality improves" remains hypothetical.

Acceptance is the necessary precondition for quality improvement—without user engagement, quality gains are impossible. Our gate-focused approach validated the critical go/no-go decision (will users engage?) while deferring comprehensive outcome evaluation to production deployment where A/B testing, expert raters, and timing instrumentation are feasible.

\textbf{Limitation 2: Self-selected early adopters.} Volunteers who opt into pilot studies are systematically different from mandatory-adoption populations. True population-level acceptance may be 10-20 percentage points lower than observed 92\%. The 22-percentage-point margin above our 70\% threshold provides buffer for selection bias.

\textbf{Limitation 3: Short-term deployment.} We measured acceptance during initial 2-week exposure. Novelty effects could inflate early acceptance, or conversely, familiarity could increase acceptance as users learn the system. Long-term trends remain unknown.

\textbf{Limitation 4: Adversarial resistance untested.} We did not evaluate how the system performs when users are incentivized to minimize effort. Users might accept low-quality suggestions to complete documentation quickly without engaging thoughtfully.

\textbf{Limitation 5: Representative corpus structure used in proof-of-concept.} Our design describes a production architecture using 500+ high-quality documentation examples for few-shot prompting. For this proof-of-concept, we used a representative corpus structure to validate the few-shot prompting mechanism. Production deployment quality will depend on systematic corpus curation effort.

\subsubsection{Broader Impact}

Improved dataset documentation benefits the ML research ecosystem by enabling accurate reproduction of findings, facilitating identification of dataset biases and limitations, supporting responsible data use, and reducing technical debt from undocumented assumptions. However, realizing these benefits requires validating that engagement translates to quality improvements.

Over-reliance on AI suggestions could reduce critical thinking about documentation content—researchers might accept suggestions without verifying accuracy or skip sections where suggestions are unavailable. If suggestions are biased (emphasizing certain use cases, reflecting English-dominant training), they could propagate systematic blind spots in documentation practices.

Researchers with limited English proficiency may face higher friction using a system trained predominantly on English exemplars, potentially exacerbating existing advantages for native English speakers if the system becomes mandatory.
